\documentclass[tikz,border=4pt]{standalone}
\usepackage{ctex}
\usepackage{amsmath}
\usepackage{pgfplots}
\pgfplotsset{compat=1.18}
\usetikzlibrary{calc, arrows.meta}

\begin{document}
% part14/05 新定义题型示意图（以自定义运算为例）
% 示意：图象对称性，展示 f(x) 与 f(-x) 的对称关系
\begin{tikzpicture}
\begin{axis}[
  axis lines=center,
  xlabel={$x$}, ylabel={$y$},
  every axis x label/.style={at={(axis cs:4,0)}, anchor=west},
  every axis y label/.style={at={(axis cs:0,7)}, anchor=south},
  xmin=-3.5, xmax=4, ymin=-1, ymax=7,
  xtick={-3,-2,-1,0,1,2,3},
  ytick={0,1,2,3,4,5,6},
  tick label style={font=\small},
  width=7.5cm, height=7cm,
  thick,
]
  % 原函数 y = x^2 + x (不对称，便于显示差别)
  \addplot[blue, thick, domain=-3:2.5, samples=80] {x^2+x};
  \node[blue, right, font=\small] at (axis cs:2,6) {$f(x)=x^2+x$};

  % 翻转函数 y = x^2 - x (将x替换为-x)
  \addplot[red, dashed, thick, domain=-2.5:3, samples=80] {x^2-x};
  \node[red, right, font=\small] at (axis cs:2.5,4) {$f(-x)=x^2-x$};

  % 对称轴 x=0
  \draw[gray, dashed] (axis cs:0,-0.5) -- (axis cs:0,7);
  \node[gray, above, font=\small] at (axis cs:0.1,7) {$y$轴对称};

  % 示意点
  \fill[blue] (axis cs:2,6) circle(2pt);
  \fill[red] (axis cs:-2,6) circle(2pt);
  \draw[gray, thin, <->] (axis cs:-2,6.3) -- (axis cs:2,6.3)
    node[midway, above, gray, font=\footnotesize] {关于$y$轴对称};

  \node[below left, font=\small] at (axis cs:0,0) {$O$};
\end{axis}
\end{tikzpicture}
\end{document}
